\documentclass[11pt]{article}
\usepackage[margin=0.88in]{geometry}
\usepackage{amsmath,amssymb,amsthm,mathtools}
\usepackage{graphicx}
\usepackage{microtype}
\usepackage[hidelinks]{hyperref}

\newtheorem{theorem}{Theorem}
\newtheorem{proposition}[theorem]{Proposition}
\newtheorem{lemma}[theorem]{Lemma}
\newtheorem{remark}[theorem]{Remark}
\newcommand{\C}{\mathbb C}
\newcommand{\R}{\mathbb R}
\newcommand{\N}{\mathbb N}
\newcommand{\tr}{\operatorname{tr}}
\newcommand{\capc}{\operatorname{cap}}
\newcommand{\D}{\Delta}

\title{A Phase-Triangular Case of the Block-GUE\\Fuglede--Kadison Determinant Conjecture}
\author{FA\_BANACH\_001, agent\_lane\_04}
\date{11 August 2026}

\begin{document}
\maketitle

\begin{abstract}
Mai and Speicher conjecture convergence of normalized determinants of arbitrary
matrix-coefficient GUE pencils to the Fuglede--Kadison determinant of the free
semicircular pencil.  We prove the full conclusion for pencils which become
upper triangular under fixed left/right equivalence and whose diagonal scalar
pencils are real up to individual phases.  This includes all real scalar pencils,
all real simultaneously triangularizable pencils, and non-diagonal triangular
systems.  For arbitrary coefficients we also prove convergence of every fixed
regularized determinant, isolating the remaining issue as uniform logarithmic
integrability at the hard edge.  The general conjecture remains open.
\end{abstract}

\section{Source conjecture and normalization}

Mai and Speicher's Conjecture~1 on PDF page 17 of \cite{MS} is reproduced
below.

\begin{center}
\includegraphics[width=0.89\textwidth]{figures/open_problem_crop.png}
\end{center}

Let $X_1^{(N)},\ldots,X_n^{(N)}$ be independent normalized GUE matrices, so
that each empirical eigenvalue law tends to the standard semicircle law on
$[-2,2]$.  For an $r\times r$ matrix $T$, write
\[
  \D_r(T)=|\det T|^{1/r}.
\]
For $a_i\in M_m(\C)$ put
\[
 A_N=\sum_{i=1}^n a_i\otimes X_i^{(N)},\qquad
 \eta_a(c)=\sum_{i=1}^n a_i c a_i^*,
\]
and
\[
 \capc(\eta_a)=\inf\{\det\eta_a(c):c>0,\ \det c=1\}.
\]
The source proves for $S_a=\sum_i a_i\otimes s_i$, with $(s_i)$ a free
standard semicircular family, that
\begin{equation}\label{eq:source-formula}
 \D(S_a)=\capc(\eta_a)^{1/(2m)}e^{-1/2}.
\end{equation}

\paragraph{Definition (phase-triangular).}
The tuple $(a_i)_{i=1}^n$ is \emph{phase-triangular} if there are
$P,Q\in GL_m(\C)$ for which $b_i=Pa_iQ$ is upper triangular for every $i$,
and, for every $1\leq j\leq m$, there are $\theta_j\in\R$ and
$r_{1j},\ldots,r_{nj}\in\R$ such that
\begin{equation}\label{eq:phase-real}
 (b_i)_{jj}=e^{\mathrm i\theta_j}r_{ij}\qquad(1\leq i\leq n).
\end{equation}
Real simultaneously upper-triangular tuples are examples (take $P=I$ and
$Q$ to be the inverse change of basis).  Allowing unrelated $P,Q$ is a
strict-equivalence version of the same flag structure.

\section{Idea of the proof}

Fixed left/right equivalence commutes with evaluation of the pencil and changes
the normalized determinant by an explicit constant.  After this equivalence,
$A_N$ is block upper triangular.  Its determinant is therefore exactly the
product of the determinants of its $N\times N$ diagonal blocks.  Condition
\eqref{eq:phase-real} makes each diagonal block a phase times a scalar multiple
of a standard GUE matrix.  The problem has thus become finitely many copies of
the classical GUE log-determinant law; no independence between these copies is
needed.

It remains to identify the resulting product with capacity.  Apply the source
formula \eqref{eq:source-formula} to the triangular free semicircular pencil.
The Fuglede--Kadison determinant of an upper triangular operator matrix is the
geometric mean of the determinants of its diagonal entries.  Each diagonal
entry is a scalar semicircular variable, whose determinant is $e^{-1/2}$ times
its standard deviation.  This gives the capacity product exactly.

\section{The phase-triangular theorem}

We use two standard determinant facts.  First, if $Z_N$ is normalized GUE,
then the GUE log-determinant asymptotics (for example, the stronger fluctuation
theorem of Bourgade--Mody \cite{BM}) imply
\begin{equation}\label{eq:GUEdet}
  \D_N(Z_N)\longrightarrow e^{-1/2}
  \quad\text{almost surely and in }L^1.
\end{equation}
Indeed, $N^{-1}\mathbb E\log|\det Z_N|\to-1/2$ and
$\operatorname{Var}(\log|\det Z_N|)=O(\log N)$; Chebyshev and
Borel--Cantelli yield the almost-sure assertion.  Uniform integrability of
$\D_N(Z_N)$ follows from
\[
 \D_N(Z_N)^2\leq \frac1N\operatorname{Tr}(Z_NZ_N^*)
\]
and the uniform Gaussian moment bounds.

Second, for an upper triangular $T=(t_{jk})\in M_m(\mathcal M)$ whose
diagonal entries are injective and of determinant class,
\begin{equation}\label{eq:triFK}
 \D_{M_m(\mathcal M)}(T)
   =\left(\prod_{j=1}^m\D_{\mathcal M}(t_{jj})\right)^{1/m}.
\end{equation}
This is the triangular Fuglede--Kadison determinant identity; see the
determinant-class discussion in \cite{Luck}.  One may also see it directly in
the algebra of affiliated operators: factor out the diagonal matrix, leaving
a unit upper-triangular matrix.  The latter is a product of elementary
unipotents and has determinant one; extended determinant multiplicativity
gives \eqref{eq:triFK}.

\begin{theorem}\label{thm:triangular}
Let $a_1,\ldots,a_n\in M_m(\C)$ be phase-triangular, and suppose that
$\eta_a$ is rank non-decreasing.  Then $A_N$ is invertible almost surely for
every $N$, and
\begin{equation}\label{eq:mainlimit}
 \D_{mN}(A_N)\longrightarrow
 \capc(\eta_a)^{1/(2m)}e^{-1/2}
\end{equation}
almost surely and in $L^1$.  Consequently the convergence also holds in
expectation.  In particular this proves the full conclusion of Conjecture~1
for every phase-triangular integer tuple.
\end{theorem}

\begin{proof}
Choose $P,Q$ and $b_i=Pa_iQ$ as in the definition, and set
\[
 B_N=\sum_i b_i\otimes X_i^{(N)}
     =(P\otimes I_N)A_N(Q\otimes I_N).
\]
The matrix $B_N$ is block upper triangular.  Its $j$th diagonal block is
\[
 Y_{j,N}=\sum_i(b_i)_{jj}X_i^{(N)}
   =e^{\mathrm i\theta_j}\sqrt{v_j}\,Z_{j,N},
 \qquad v_j=\sum_i r_{ij}^2,
\]
where $Z_{j,N}=\sum_i r_{ij}X_i^{(N)}/\sqrt{v_j}$ is marginally standard
GUE.  Every $v_j$ is positive: otherwise the triangular free pencil
$\sum_i b_i\otimes x_i$ has a zero diagonal entry and cannot have full free
rank, contrary to the rank-non-decrease criterion recalled in \cite{MS}.
Thus each $Y_{j,N}$, and hence $B_N$ and $A_N$, is almost surely invertible.

Finite block-triangular factorization gives
\begin{equation}\label{eq:finitefactor}
 \D_{mN}(B_N)=\left(\prod_{j=1}^m\D_N(Y_{j,N})\right)^{1/m}.
\end{equation}
By \eqref{eq:GUEdet}, simultaneously for the finitely many $j$,
\[
 \D_{mN}(B_N)\longrightarrow
 e^{-1/2}\left(\prod_{j=1}^m v_j\right)^{1/(2m)}
\]
almost surely.  The same convergence holds in $L^1$: the geometric mean in
\eqref{eq:finitefactor} is at most the arithmetic mean, and the proof of
\eqref{eq:GUEdet} gives uniform $L^2$ bounds for its finitely many factors.
Moreover,
\begin{equation}\label{eq:PQfinite}
 \D_{mN}(B_N)=|\det P\det Q|^{1/m}\D_{mN}(A_N).
\end{equation}

We now identify the constant.  Let $S_b=\sum_i b_i\otimes s_i$.  It is upper
triangular with diagonal entries
$e^{\mathrm i\theta_j}\sqrt{v_j}\,s'_j$, where each $s'_j$ is a standard
semicircular variable.  Since
$\D(s'_j)=\exp\int\log|t|\,d\mu_{\rm sc}(t)=e^{-1/2}$,
\eqref{eq:triFK} and the source formula give
\[
 \capc(\eta_b)^{1/(2m)}e^{-1/2}
 =\D(S_b)
 =e^{-1/2}\left(\prod_jv_j\right)^{1/(2m)}.
\]
Hence $\capc(\eta_b)=\prod_jv_j$.

Finally, directly from the capacity definition,
\begin{equation}\label{eq:capcov}
 \capc(\eta_b)=|\det P\det Q|^2\capc(\eta_a).
\end{equation}
Indeed $\eta_b(c)=P\eta_a(QcQ^*)P^*$, and
$|\det Q|^{-2/m}QcQ^*$ runs through all positive determinant-one matrices as
$c$ does.  Combining \eqref{eq:PQfinite}--\eqref{eq:capcov} proves
\eqref{eq:mainlimit}.
\end{proof}

\section{Two general reductions}

The first reduction gives genuine invertibility progress without a triangular
hypothesis.  Let $\mathcal A=\operatorname{span}\{a_1,\ldots,a_n\}$.

\begin{proposition}[blow-up subsequence]\label{prop:blowup}
Suppose $\eta_a$ is rank non-decreasing.  There is a positive integer $d$
(one may take $d\leq m-1$ when $m>1$ by \cite{DM}) such that, for every
$k\geq1$, $A_{kd}$ is invertible almost surely.
\end{proposition}

\begin{proof}
Rank non-decrease is equivalent to full noncommutative rank.  Thus some
$d$th blow-up of $\mathcal A$ contains an invertible matrix, and direct sums
give an invertible evaluation in every blow-up of size $kd$.  Consequently
the polynomial
$\det(\sum_i a_i\otimes Z_i)$ on $M_{kd}(\C)^n$ is nonzero.  Hermitian tuples
form a real form, hence are Zariski dense in the complex tuple space, so the
restriction of this polynomial to Hermitian tuples is still nonzero.  Its zero
set has Lebesgue measure zero there, while a GUE tuple has a smooth positive
density.  Therefore $A_{kd}$ is almost surely invertible.
\end{proof}

The second reduction holds in every size and identifies the analytic gap.
For $\varepsilon>0$ set
\[
 \D_{N,\varepsilon}(A_N)
 =\det(A_NA_N^*+\varepsilon^2I_{mN})^{1/(2mN)}.
\]

\begin{proposition}[Regularized general reduction]\label{prop:regularized}
For arbitrary fixed $a_i\in M_m(\C)$ and every $\varepsilon>0$,
\begin{equation}\label{eq:reglimit}
 \D_{N,\varepsilon}(A_N)\longrightarrow
 \exp\!\left\{\frac12(\tr_m\otimes\tau)
 \log(S_aS_a^*+\varepsilon^2I)\right\}
\end{equation}
almost surely and in $L^1$.  If $\eta_a$ is rank non-decreasing, the
right-hand side decreases, as $\varepsilon\downarrow0$, to
$\capc(\eta_a)^{1/(2m)}e^{-1/2}$.
\end{proposition}

\begin{proof}
Almost-sure asymptotic freeness of independent GUE matrices gives convergence
of the empirical spectral law of $A_NA_N^*$ to that of $S_aS_a^*$, together
with an almost-sure common compact upper bound on the spectra.  For fixed
$\varepsilon$, the function $t\mapsto\frac12\log(t+\varepsilon^2)$ is
continuous on this compact interval, proving \eqref{eq:reglimit}.  The
arithmetic--geometric mean bound and uniform Gaussian operator/Hilbert--Schmidt
moment bounds give uniform integrability, hence $L^1$ convergence.  Monotone
convergence in the logarithmic functional calculus gives the
$\varepsilon\downarrow0$ limit, and \eqref{eq:source-formula} identifies it.
\end{proof}

\section{Hard-edge obstruction and scope}

If $s_1(A_N),\ldots,s_{mN}(A_N)$ are the singular values, the difference
between regularized and actual logarithmic determinants is exactly
\[
 \log\D_{N,\varepsilon}(A_N)-\log\D_{mN}(A_N)
 =\frac1{2mN}\sum_{k=1}^{mN}
   \log\!\left(1+\frac{\varepsilon^2}{s_k(A_N)^2}\right).
\]
Thus Proposition~\ref{prop:regularized} reduces the full conjecture to showing
that this nonnegative quantity vanishes uniformly in the large-$N$ limit as
$\varepsilon\downarrow0$, almost surely and in mean.  This is precisely a
hard-edge uniform-integrability estimate.  Strong convergence controls the
top edge, while generic polynomial anti-concentration loses a factor growing
with the determinant degree $mN$.  Existing matrix-Dyson local laws for
Kronecker ensembles require stability or flatness hypotheses not verified here
for every merely rank non-decreasing covariance map.

Accordingly, Theorem~\ref{thm:triangular} is a solved structural subcase and
Propositions~\ref{prop:blowup}--\ref{prop:regularized} are general reductions;
they do not settle arbitrary coefficient tuples.  Bounded run-index and web
searches found no exact later solution or duplicate of the phase-triangular
statement.  The novelty claim remains provisional pending specialist review.

\begin{thebibliography}{9}
\bibitem{BM}
P. Bourgade and K. Mody,
\emph{Gaussian fluctuations of the determinant of Wigner matrices},
Electron. J. Probab. 24 (2019), Paper No. 96; arXiv:1811.06815.

\bibitem{DM}
H. Derksen and V. Makam,
\emph{On non-commutative rank and tensor rank},
Linear Multilinear Algebra 66 (2018), 1069--1084; arXiv:1606.06701.

\bibitem{Luck}
W. L\"uck,
\emph{$L^2$-invariants: Theory and applications to geometry and $K$-theory},
Ergebnisse der Mathematik, vol. 44, Springer, 2002.

\bibitem{MS}
T. Mai and R. Speicher,
\emph{Fuglede--Kadison determinants of matrix-valued semicircular elements and
capacity estimates}, arXiv:2406.15922 (2024).
\end{thebibliography}

\end{document}
